% Блок-схемы parse_factor / parse_function / apply_function / parse_number
% ГОСТ 19.701-90, src/parser.cpp

\documentclass[a4paper]{article}
\usepackage[T2A]{fontenc}
\usepackage[utf8]{inputenc}
\usepackage[russian]{babel}
\usepackage[a4paper,margin=1.4cm]{geometry}
\usepackage{tikz}
\usetikzlibrary{shapes.geometric, arrows.meta, positioning, calc}

\tikzset{
  terminal/.style={
    rectangle, rounded corners=8pt,
    draw=black, line width=1pt,
    minimum width=4cm, minimum height=0.85cm,
    text centered, font=\footnotesize
  },
  process/.style={
    rectangle,
    draw=black, line width=1pt,
    minimum width=4cm, minimum height=0.85cm,
    text centered, font=\footnotesize
  },
  decision/.style={
    diamond, aspect=2.3,
    draw=black, line width=1pt,
    inner sep=1pt,
    minimum width=3cm, minimum height=0.85cm,
    text centered, font=\footnotesize
  },
  predefined/.style={
    rectangle, double, double distance=2pt,
    draw=black, line width=1pt,
    minimum width=4cm, minimum height=0.85cm,
    text centered, font=\footnotesize
  },
  connector/.style={
    circle,
    draw=black, line width=1pt,
    minimum size=0.65cm,
    text centered, font=\footnotesize\bfseries
  },
  arrow/.style={->, >=Stealth, line width=1pt}
}

% ================================================================
% ЧАСТЬ 1 — parse_factor()
%
% Принцип: каждая ветка «Да» = один блок СТРОГО НА УРОВНЕ ромба.
% Никаких вертикально накладывающихся правых ветвей.
% Расстояния между ромбами подобраны с учётом высоты правых блоков.
% ================================================================
\begin{document}
\begin{center}
\textbf{Блок-схема функции \texttt{parse\_factor()}}
\end{center}

\begin{center}
\begin{tikzpicture}[node distance=0.9cm]

\node (start)  [terminal]                        {Начало parse\_factor};
\node (getC)   [predefined, below=of start]      {c = peek()};

\node (isMin)  [decision,   below=of getC]       {c == \textquoteleft-\textquoteright?};
\node (isPar)  [decision,   below=1.8cm of isMin]{c == \textquoteleft(\textquoteright?};
\node (isAlph) [decision,   below=1.6cm of isPar]{isalpha(c)?};
\node (isDig)  [decision,   below=1.6cm of isAlph]{isdigit(c) или c==\textquoteleft.\textquoteright?};
\node (isEos)  [decision,   below=1.6cm of isDig]{c == \textquoteleft{\textbackslash}0\textquoteright?};
\node (errSym) [terminal,   below=1.2cm of isEos]{throw \textquotedbl Неожиданный символ\textquotedbl};

%% Правые блоки — каждый на уровне своего ромба (right=4.5cm)
\node (doMin)  [predefined, right=4.5cm of isMin]
               {eat(); \ return $-$parse\_factor()};

\node (doPar)  [predefined, right=4.5cm of isPar, text width=5.5cm, align=center]
               {eat(); r = parse\_expression();\\expect(\textquoteleft)\textquoteright); \ return r};

\node (doAlph) [predefined, right=4.5cm of isAlph]
               {return parse\_function()};

\node (doDig)  [predefined, right=4.5cm of isDig]
               {return parse\_number()};

\node (errEos) [terminal,   right=4.5cm of isEos]
               {throw \textquotedbl Неожиданный конец\textquotedbl};

%% Стрелки главной оси
\draw [arrow] (start)  -- (getC);
\draw [arrow] (getC)   -- (isMin);
\draw [arrow] (isMin)  -- node[left]{\footnotesize Нет} (isPar);
\draw [arrow] (isPar)  -- node[left]{\footnotesize Нет} (isAlph);
\draw [arrow] (isAlph) -- node[left]{\footnotesize Нет} (isDig);
\draw [arrow] (isDig)  -- node[left]{\footnotesize Нет} (isEos);
\draw [arrow] (isEos)  -- node[left]{\footnotesize Нет} (errSym);

%% Стрелки правых веток (Да)
\draw [arrow] (isMin)  -- node[above]{\footnotesize Да} (doMin);
\draw [arrow] (isPar)  -- node[above]{\footnotesize Да} (doPar);
\draw [arrow] (isAlph) -- node[above]{\footnotesize Да} (doAlph);
\draw [arrow] (isDig)  -- node[above]{\footnotesize Да} (doDig);
\draw [arrow] (isEos)  -- node[above]{\footnotesize Да} (errEos);

\end{tikzpicture}
\end{center}

\newpage

% ================================================================
% ЧАСТЬ 2 — parse_function()
%
% Цикл чтения имени. Возврат цикла — по правому полю:
% addCh → вправо → вверх до cond → влево к cond.east
% (Z-путь: три сегмента, x одинаковый на вертикальном участке)
% ================================================================
\begin{center}
\textbf{Блок-схема функции \texttt{parse\_function()}}
\end{center}

\begin{center}
\begin{tikzpicture}[node distance=0.8cm]

\node (start)  [terminal]                    {Начало parse\_function};
\node (init)   [process,    below=of start]  {name = \textquotedbl\textquotedbl};
\node (cond)   [decision,   below=of init]   {isalpha(s[pos])?};
\node (expL)   [process,    below=of cond]   {expect(\textquoteleft(\textquoteright)};
\node (arg)    [predefined, below=of expL]   {arg = parse\_expression()};
\node (expR)   [process,    below=of arg]    {expect(\textquoteleft)\textquoteright)};
\node (apply)  [predefined, below=of expR]   {r = apply\_function(name, arg)};
\node (ret)    [terminal,   below=of apply]  {Возврат r};

\node (addCh)  [process,    right=4cm of cond]{name += s[pos]; \ pos++};

\draw [arrow] (start)  -- (init);
\draw [arrow] (init)   -- (cond);
\draw [arrow] (cond)   -- node[above]{\footnotesize Да}  (addCh);
\draw [arrow] (cond)   -- node[left] {\footnotesize Нет} (expL);
\draw [arrow] (expL)   -- (arg);
\draw [arrow] (arg)    -- (expR);
\draw [arrow] (expR)   -- (apply);
\draw [arrow] (apply)  -- (ret);

%% Возврат addCh → cond: «мостик» над cond.
%% addCh стоит на том же уровне что cond, поэтому нельзя идти «вверх до cond».
%% Вместо этого: вправо → вверх выше cond.north → влево+вниз к cond.north.
%% Используем -| (сначала вертикаль, потом горизонталь), входим в cond сверху.
\draw [arrow] (addCh.east) -- ++(0.8,0) -- ++(0,0.7) -| (cond.north);

\end{tikzpicture}
\end{center}

\newpage

% ================================================================
% ЧАСТЬ 3 — apply_function(name, x)
%
% Каскад ромбов. Для каждой функции — правая подсхема.
% Расстояние 3.2cm между ромбами: правые блоки не перекрываются.
% ================================================================
\begin{center}
\textbf{Блок-схема функции \texttt{apply\_function(name,\ x)}}
\end{center}

\begin{center}
\begin{tikzpicture}[node distance=0.8cm]

\node (start)  [terminal]                         {Начало apply\_function};
\node (isLn)   [decision,   below=of start]       {name == \textquotedbl ln\textquotedbl?};
\node (isExp)  [decision,   below=3.4cm of isLn]  {name == \textquotedbl exp\textquotedbl?};
\node (isSqrt) [decision,   below=3.4cm of isExp] {name == \textquotedbl sqrt\textquotedbl?};
\node (isAbs)  [decision,   below=3.4cm of isSqrt]{name == \textquotedbl abs\textquotedbl?};
\node (errUnk) [terminal,   below=of isAbs]       {throw \textquotedbl Неизвестная функция\textquotedbl};

%% Ветка ln
\node (chkLn)  [decision,   right=4.5cm of isLn]  {$x \leq 0$?};
\node (errLn)  [terminal,   above=0.7cm of chkLn] {throw \textquotedbl ln: $x \leq 0$\textquotedbl};
\node (retLn)  [terminal,   below=0.7cm of chkLn] {Возврат $\ln(x)$};

%% Ветка exp
\node (calcE)  [process,    right=4.5cm of isExp] {r = exp($x$)};
\node (chkFin) [decision,   below=0.7cm of calcE] {!isfinite(r)?};
\node (errOvf) [terminal,   right=2.5cm of chkFin]{throw \textquotedbl Переполнение\textquotedbl};
\node (retExp) [terminal,   below=0.7cm of chkFin]{Возврат r};

%% Ветка sqrt
\node (chkSq)  [decision,   right=4.5cm of isSqrt]{$x < 0$?};
\node (errSq)  [terminal,   above=0.7cm of chkSq] {throw \textquotedbl sqrt: $x < 0$\textquotedbl};
\node (retSq)  [terminal,   below=0.7cm of chkSq] {Возврат $\sqrt{x}$};

%% Ветка abs
\node (retAbs) [terminal,   right=4.5cm of isAbs] {Возврат $|x|$};

%% Главная ось
\draw [arrow] (start)  -- (isLn);
\draw [arrow] (isLn)   -- node[left]{\footnotesize Нет} (isExp);
\draw [arrow] (isExp)  -- node[left]{\footnotesize Нет} (isSqrt);
\draw [arrow] (isSqrt) -- node[left]{\footnotesize Нет} (isAbs);
\draw [arrow] (isAbs)  -- node[left]{\footnotesize Нет} (errUnk);

%% Правые ветки
\draw [arrow] (isLn)   -- node[above]{\footnotesize Да} (chkLn);
\draw [arrow] (chkLn)  -- node[left] {\footnotesize Да} (errLn);
\draw [arrow] (chkLn)  -- node[left] {\footnotesize Нет}(retLn);

\draw [arrow] (isExp)  -- node[above]{\footnotesize Да} (calcE);
\draw [arrow] (calcE)  -- (chkFin);
\draw [arrow] (chkFin) -- node[above]{\footnotesize Да} (errOvf);
\draw [arrow] (chkFin) -- node[left] {\footnotesize Нет}(retExp);

\draw [arrow] (isSqrt) -- node[above]{\footnotesize Да} (chkSq);
\draw [arrow] (chkSq)  -- node[left] {\footnotesize Да} (errSq);
\draw [arrow] (chkSq)  -- node[left] {\footnotesize Нет}(retSq);

\draw [arrow] (isAbs)  -- node[above]{\footnotesize Да} (retAbs);

\end{tikzpicture}
\end{center}

\newpage

% ================================================================
% ЧАСТЬ 4 — parse_number()
%
% Вершина цикла — conn[3].
% condRd: Нет → вниз (пост-цикл), Да → вправо (тело цикла).
% Тело: isDot → Нет (вниз) → incPos; Да (вправо) → chkHD.
% Возвраты цикла — по правому полю, каждый на своём X:
%   incPos: X = isDot.east.x + 1.5
%   setHD:  X = chkHD.east.x + 1.5
% Оба идут вверх до conn[3] и входят в него сверху.
% ================================================================
\begin{center}
\textbf{Блок-схема функции \texttt{parse\_number()}}
\end{center}

\begin{center}
\begin{tikzpicture}[node distance=0.8cm]

%% ---- Главная ось ----
\node (start)   [terminal]                      {Начало parse\_number};
\node (init)    [process,    below=of start]    {start = pos; \ has\_dot = false};
\node (conn3)   [connector,  below=of init]     {3};           %% вершина цикла
\node (condRd)  [decision,   below=of conn3]
                {isdigit(s[pos]) или s[pos]==\textquoteleft.\textquoteright?};

%% Пост-цикл (Нет — вниз)
\node (mkNum)   [process,    below=of condRd]   {num = s.substr(start,\ pos$-$start)};
\node (chkDot)  [decision,   below=of mkNum]    {num == \textquotedbl.\textquotedbl?};
\node (errDot)  [terminal,   right=4cm of chkDot]{throw \textquotedbl Одиночная точка\textquotedbl};
\node (intLoop) [process,    below=of chkDot]   {Цикл: целая часть $\to$ value};
\node (hasFrac) [decision,   below=of intLoop]  {num[k] == \textquotedbl.\textquotedbl?};
\node (fracLp)  [process,    right=4cm of hasFrac]{Цикл: дробная часть $\to$ value};
\node (ret)     [terminal,   below=of hasFrac]  {Возврат value};

%% ---- Тело цикла (Да — вправо от condRd) ----
\node (isDot)   [decision,   right=5cm of condRd]{s[pos] == \textquoteleft.\textquoteright?};
\node (incPos)  [process,    below=1.2cm of isDot]{pos++};     %% Нет-ветка isDot
\node (chkHD)   [decision,   right=3.5cm of isDot]{has\_dot?}; %% Да-ветка isDot
\node (err2dot) [terminal,   above=0.7cm of chkHD]{throw \textquotedbl Две точки\textquotedbl};
\node (setHD)   [process,    below=0.7cm of chkHD]{has\_dot=true; \ pos++};

%% ---- Стрелки главной оси ----
\draw [arrow] (start)   -- (init);
\draw [arrow] (init)    -- (conn3);
\draw [arrow] (conn3)   -- (condRd);
\draw [arrow] (condRd)  -- node[left] {\footnotesize Нет}(mkNum);
\draw [arrow] (mkNum)   -- (chkDot);
\draw [arrow] (chkDot)  -- node[above]{\footnotesize Да} (errDot);
\draw [arrow] (chkDot)  -- node[left] {\footnotesize Нет}(intLoop);
\draw [arrow] (intLoop) -- (hasFrac);
\draw [arrow] (hasFrac) -- node[above]{\footnotesize Да} (fracLp);
\draw [arrow] (hasFrac) -- node[left] {\footnotesize Нет}(ret);
\draw [arrow] (fracLp.south) -- ++(0,-0.4) -| (ret.east);

%% ---- Стрелки тела цикла ----
\draw [arrow] (condRd)  -- node[above]{\footnotesize Да} (isDot);
\draw [arrow] (isDot)   -- node[left] {\footnotesize Нет}(incPos);
\draw [arrow] (isDot)   -- node[above]{\footnotesize Да} (chkHD);
\draw [arrow] (chkHD)   -- node[left] {\footnotesize Да} (err2dot);
\draw [arrow] (chkHD)   -- node[left] {\footnotesize Нет}(setHD);

%% ---- Возврат incPos → conn3 ----
%% incPos.east ≈ x=7; chkHD.east ≈ x=10, setHD.east ≈ x=10.5
%% Нужен x > 10.5, чтобы вертикаль не проходила сквозь chkHD/setHD.
%% Выбираем x = incPos.east + 4 ≈ 11. Входим в conn3 сверху (north).
\coordinate (IbotR) at ($(incPos.east) + (4.0, 0)$);
\coordinate (ItopR) at ($(IbotR        |- conn3.north)$);
\draw [arrow] (incPos.east) -- (IbotR) -- (ItopR) -- (conn3.north);

%% ---- Возврат setHD → conn3 ----
%% setHD.east ≈ x=10.5; используем x = setHD.east + 2 ≈ 12.5 (другой X).
%% Входим в conn3 сбоку (east) — другая точка, нет наложения с путём incPos.
\coordinate (SbotR) at ($(setHD.east) + (2.0, 0)$);
\coordinate (StopR) at ($(SbotR       |- conn3.east)$);
\draw [arrow] (setHD.east) -- (SbotR) -- (StopR) -- (conn3.east);

\end{tikzpicture}
\end{center}

\end{document}
